\begingroup
\newcommand{\ctP}{\mathsf P}
\newcommand{\ctI}{\mathsf I}
\newcommand{\ctQ}{\mathsf Q}
\newcommand{\ctD}{\mathcal D}
\newcommand{\ctS}{\mathsf S}
\newcommand{\ctMP}{\mathsf{MP}}
\newcommand{\ctAX}{\mathsf{AX}}
\newcommand{\ctH}{\mathsf H}
\newcommand{\ctPush}{\mathsf{PUSH}}
\newcommand{\ctDiscard}{\mathsf{DISCARD}}
\newcommand{\ctBlock}{\mathcal B}

\section{A 201-subtree source: transfer-pop and cache-return inference}
\label{sec:cache-calculus}

This construction combines the transfer-pop proof enumerator
of~\cite{transfer} with cache-return modus ponens from~\cite{cachemp},
rebased onto the accepted K215 source~\cite{k215}. The combination needs
a new proof compiler: cache-return inference no longer supplies the
dummy stack entry consumed when its result is transferred to the current
formula. Reserving that entry before compiling the premises resolves
the mismatch. The resulting literal programs have 201 distinct source
subtrees, one fewer than the accepted source. This section proves their
ordinary source-level consistency calibration. It does not identify a
214-state binary machine or infer a physical state count from the
source count; allocation, full-tape simulation, and blank initialization
are separate obligations.

\subsection{Numerical calculus and exact command interface}

Use Cantor pairing and the unchanged implication and quantifier codes:
\[
 \ctP(a,b)=\frac{(a+b)(a+b+1)}2+b,\qquad
 \ctI(a,b)=4\ctP(a,b),\qquad
 \ctQ_i a=4\ctP(i,a)+2.
\]
Truth is code0 and falsity is code1. In particular,
$\ctI(0,0)=0$. Write $[x;w]=\ctP(x,w)$ and use right-nested
pairing for longer stacks. This is numerical notation, not an
injective encoding of finite lists: $[0;0]=0$, and both projections
of zero are zero.

Let $\ctD$ be the accepted dominated numerical calculus: truth,
numerical MP, theorem-only generalization, the seven retained logical
rows $R_1,\ldots,R_7$, the complete gated theory bundle, and
\[
 L(A,B,C,U)=
 \ctI\bigl(\ctI(\ctI(C,U),A),
             \ctI(\ctI(A,C),\ctI(B,C))\bigr),\qquad U\leq A.
\]
These rows and the full bundle are retained literally. In particular,
none of the bundle's support or atomic-vacuity material is removed.
The inherited finite propositional certificate recovers unrestricted
$\rho_8$ for every natural code; the accepted calibration then gives
\begin{equation}
 \ctD\vdash1\quad\Longleftrightarrow\quad
 \exists p\;\operatorname{Proof}_{\mathrm{ZFC}}(p,\bot).
 \label{eq:ct-inherited}
\end{equation}
The present theorem uses this existing mathematical
calibration~\cite{k215}; it is not a new proof-assistant verification
of that dependency. Formula decoding and well-formedness tests are
not added to the numerical calculus.

The proof enumerator has a new operational language with unchanged
axiom strength. In particular, the operand of MP remains an
arbitrary natural increment; inference is not restricted to a fixed
consequent or a fixed parameter. Successful inference, cache movement,
and stack restoration are reorganized together. Treating either command
change in isolation would miss the obstruction solved by the reserved
zero, and would not justify transferring the old completeness theorem.

At command boundaries the work tuple is $(t,w,d,h)$: current formula,
paired stack, additive parameter, and cached formula. If
$\operatorname{unpair}(w)=(x,v)$, the direct transfer command is
\[
 \ctS(t,w,d,h)=
 \begin{cases}
  \mathsf{halt},&h=1,\\
  (h,v,d,x),&h\ne1.
 \end{cases}
\]
For $Y=d+x$, cache-return inference has the exact equation
\[
 \ctMP(x)(t,w,d,h)=
 \begin{cases}
  (0,w,0,Y),&\ctI(t,Y)=h,\\
  (0,w,Y,0),&\ctI(t,Y)\ne h.
 \end{cases}
\]
The successful branch clears the parameter and leaves the conclusion
in cache; it does not push a dummy zero. The failed branch is unchanged.

The axiom command $\ctAX(c;x,y)$ is unchanged. Set
$U=d$, $A=d+x$, $B=h+y$, and $C=c$. Its result is
\begin{align*}
 t'&=G(A,B,c),& d'&=0,&h'&=0,\\
 w'&=[R_1,\ldots,R_7,L(A,B,c,U),\ctQ_U h,t;w].
\end{align*}
Here $G$ is the original full bundle when $c\ne A,B$, and zero
otherwise. The quantified entry uses the incoming cache, before adding
$y$. Every logical row is an implication code and hence differs from1;
the quantified entry is nonzero of residue2 modulo4. The gate is also
never1: an enabled bundle is implication-coded, and a disabled gate
is zero. These facts hold for arbitrary raw operands.

\subsection{Literal construction and eight variants}

In the accepted tree, replace the successful-MP subtree at
\path{11010010111} by $(\mathrm{CD}_4,
(\mathrm{CI}_5,\mathrm{CC}_2))$. Delete the DROP selector layer
by replacing \path{11010} with its current descendant
\path{11010010}. Finally replace the old NIP body at \path{111}
by a transfer body. The direct body is the closed composition
\[
 \mathsf{clearT};\mathsf{copyH};\mathsf{oldNIP}.
\]
The first part clears current; the second is the old DROP cache-copy
tail, preserving cache and clearing scratch. Old NIP then tests that
cache and pops the stack. Both parenthesizations implement precisely
the direct $\ctS$ equation.

Alternatively use
$\mathsf{oldDROP};\mathsf{oldNIP};\mathsf{oldNIP}$, abbreviated
DNN. Its nonhalting map is the same, but it also halts if incoming
$t=1$. It is therefore an early-halting refinement of direct transfer.
Each body has two parenthesizations, and each program has two small
selector orders, giving eight literal variants:
\begin{center}
\begin{tabular}{lrrrrr}
\toprule
Transfer body&Variants&Distinct subtrees&Leaves&Depth&Zero run\\
\midrule
Direct&4&201&4709&44&6\\
DNN&4&201&4752&44&6\\
\bottomrule
\end{tabular}
\end{center}
For the ordinary order, wire fields are
\[
 f(\ctS)=[0],\qquad f(\ctMP(x))=[1,x],\qquad
 f(\ctAX(c;x,y))=[c+2,y,x].
\]
The \texttt{swap01} order exchanges only the first two selectors:
$f(\ctS)=[1]$ and $f(\ctMP(x))=[0,x]$. AX fields stay fixed.
The first two selector decrements prove these partitions for all
natural fields: zero selects the first body, one the second, and
$c+2$ enters AX with residual $c$.

The independent source audit reconstructs all eight files, checks both
CFG resolvers, and verifies all 4565 primitive effects and internal
successors of AX. Its unique exterior edge still reaches the common
branch continuation. Thus AX is not incorrectly assumed closed when
its required outer continuation is supplied by the selector.

\subsection{Soundness on all raw candidate words}

Call a work tuple safe if current, cache, and every successive decoded
stack head are $\ctD$-theorems. No theoremhood condition is imposed
on $d$. AX preserves safety by its inherited axiom and generalization
contracts. Failed MP inserts truth; successful MP derives $Y$ from
$t$ and $\ctI(t,Y)$, places it in cache, and inserts truth in current.
Transfer only moves theorem codes. Its cache-one guard, or either guard
of DNN, therefore witnesses a theorem equal to falsity. Consequently
every actual halt from a safe seed yields $\ctD\vdash1$. This argument
quantifies over all natural command fields, including unintended or
nondecodable inputs; it does not presuppose a valid compiled proof word.

\subsection{Structural blocks recovered without DROP}

Put $D_0=\ctAX(0;0,0)$ and define
\begin{align*}
 \ctH&=D_0;\ctS^{11},\\
 \ctPush&=D_0;\ctS^9;\ctMP(0),\\
 \ctDiscard&=\ctMP(1);\ctAX(1;0,0);\ctS^{10};\ctMP(0).
\end{align*}
From $(t,[x;w],d,h)$, the first ten shifts of $\ctH$ yield
$(\ctQ_dh,[x;w],0,t)$. Their tests inspect only zero and the nine
generated nonfalsity entries. The eleventh shift halts exactly when
$t=1$, otherwise returning $(t,w,0,x)$. Thus $\ctH$ preserves
current while popping one old stack entry and resetting the parameter.
It intentionally does not test the original cache before AX erases it.

After nine shifts in $\ctPush$, current is the dominated row,
cache is $\ctQ_dh$, and stack is $[t;w]$. The final MP necessarily
fails because an implication code has residue0 and this cache has
residue2. Hence PUSH never halts and returns $(0,[t;w],0,0)$ for
every natural input tuple.

DISCARD starts from $(t,w,0,0)$. Its first MP fails since
$\ctI(t,1)>0$. AX has $U=A=c=1$ and gate zero. Ten shifts give
$(\ctQ_1 0,w,0,0)$; the final MP fails because
$\ctI(\ctQ_1 0,0)>0$. Thus DISCARD never halts and returns
$(0,w,0,0)$. Finally, from current, cache, and parameter zero, define
$J(0)$ to be empty and $J(u)=\ctMP(u)$ when $u>0$. The latter is
a forced failure setting $d=u$. Omitting $\ctMP(0)$ is essential:
the zero alias would make it succeed.

For DNN, the extra top test does not invalidate these exact auxiliary
claims: every tested gate and generated entry is nonfalsity. The old
top remains untested in PUSH, and its sole test in H is the intended
last shift.

These blocks use the axiom emitter as a source of temporary storage.
Its fixed ten-entry output preserves the old current formula behind nine
safe entries; repeated transfers can recover that saved value or leave
it on the stack without a primitive DROP command. The intermediate
logical rows need not be useful premises in the proof being compiled.
Their exact stack positions, their nonfalsity tags, and the quantified
entry's residue are sufficient. This distinction matters for compression:
a command that is expensive as an isolated arithmetic routine can supply
several structural operations through already shared source subtrees.
The resulting proof words can be very long. No execution-time advantage
is required for a consistency machine, whose specification asks only
whether a finite halting run exists.

\subsection{The strong theorem-block compiler}

For every derivable code $X$, construct a finite block $\ctBlock_X$
which, from \emph{every} $(t,w,0,0)$, either halts or returns exactly
$(X,w,0,0)$. Arbitrary incoming $t,w$ need not be safe. This universal
completeness contract is distinct from the safe-run soundness invariant.
Truth uses DISCARD. The axiom cases are
\begin{align*}
 R_j:\quad&\ctDiscard;\ctAX(c;a,b);\ctS^{j+1};\ctH^{9-j}
                    &&(1\leq j\leq7),\\
 L(a,b,c,u):\quad&\ctDiscard;J(u);\ctAX(c;a-u,b);\ctS^9;\ctH
                    &&(u\leq a),\\
 G(a,b,c):\quad&\ctDiscard;\ctAX(c;a,b);\ctH^{10}
                    &&(c\ne a,b).
\end{align*}
The indicated shifts select the desired row; the H repetitions preserve
it and remove exactly the remaining generated entries, ending with
saved zero in cache. The underlying stack is restored. All auxiliary
MP calls are forced failures, so cache-return success does not affect
these calculations.

Generalization uses
\[
 \ctBlock_{\ctQ_iX}=
 \ctBlock_X;\ctPush;J(i);\ctS;\ctAX(i;0,0);\ctS^{10}.
\]
After PUSH, setter, and transfer, the state is $(0,w,i,X)$. That
transfer tests incoming cache zero, not the newly popped $X$. AX has
zero gate and generates $\ctQ_iX$ immediately above saved zero;
ten shifts return $(\ctQ_iX,w,0,0)$. This includes $i=0$ and
$X=1$, subject only to permitted earlier halting.

\paragraph{MP with a protected zero.}
For premises $A$ and $\ctI(A,Y)$ define
\begin{multline*}
 \ctBlock_Y=\ctDiscard;\ctPush;
 \ctBlock_{\ctI(A,Y)};\ctPush;
 \ctBlock_A;\ctPush;\\
 \ctS;\ctS;\ctMP(Y);\ctS.
\end{multline*}
If no earlier halt occurs, its exact state trace is
\begin{align*}
 (t,w,0,0)&\longmapsto(0,[0;w],0,0)\\
 &\longmapsto(0,[\ctI(A,Y),0;w],0,0)\\
 &\longmapsto(0,[A,\ctI(A,Y),0;w],0,0)\\
 &\longmapsto(0,[\ctI(A,Y),0;w],0,A)\\
 &\longmapsto(A,[0;w],0,\ctI(A,Y))\\
 &\longmapsto(0,[0;w],0,Y)\\
 &\longmapsto(Y,w,0,0).
\end{align*}
The second transfer may halt if $A=1$, and the final transfer halts
if $Y=1$. Otherwise the displayed equations hold. Premise blocks
preserve their entire incoming stacks, which protects the reserved zero
until it compensates for the final pop. Without this reservation, that
pop would destroy an arbitrary original stack head. Pairing identities
validate the trace even at $A=Y=w=0$; no list freeness is assumed.

Induction on a finite derivation proves the strong contract. A proof
DAG can be expanded to a finite tree before compiling the major premise
first. DNN transfers the contract unchanged: runs agree with the direct
version until DNN possibly halts earlier. If $\ctBlock_1$ returns,
append H to detect its current falsity. Thus $\ctBlock_1;\ctH$
succeeds from every $(t,w,0,0)$.

\subsection{Cofinal successful candidates and the source theorem}

Let $F(fs,r)$ right-fold the command fields above residual $r$, and set
\[
 \operatorname{tail}([])=0,\qquad
 \operatorname{tail}(k::L)=1+F(f(k),\operatorname{tail}(L)).
\]
A continuation first decrements its positive proof-list pointer;
a fresh candidate decodes without that outer successor. Put
$R=\operatorname{tail}(\ctBlock_1;\ctH)$. For each natural $z$,
\[
 N_z=\ctP\bigl(2,\ctP(z,\ctP(1,R))\bigr)
\]
starts with $\ctAX(0;1,z)$ in either selector order. This resets
$d=h=0$ from every retained tuple. Its current and stack may be
arbitrary, precisely as allowed by the strong contract. Since
$N_z\geq z$, successful candidates are cofinal beyond every fixed
finite startup cutoff.

Every nonhalting positive continuation strictly decreases: after its
predecessor operation, each field reader takes a Cantor right projection
no larger than that predecessor. The natural-counter macros terminate.
Hence each nonhalting candidate drains, and sequential enumeration
reaches a successful $N_z$ unless it has already halted.

\begin{theorem}[Source consistency calibration]
For each of the eight literal transfer/cache programs, started with
the unchanged enumerator scheduler and a safe initialized frame,
halting is equivalent to $\ctD\vdash1$. Consequently, using
\eqref{eq:ct-inherited}, its nonhalting is equivalent in ZFC to
$\operatorname{Con}(\mathrm{ZFC})$.
\end{theorem}

\begin{proof}
Safe-run soundness gives the forward implication. A native refutation
gives the strong block and cofinally many successful candidates, proving
the reverse implication. These finite syntactic transformations and
natural-number arguments are formalizable in ZFC. A binary realization
requires its own exact runtime and initializer proofs.
\end{proof}

\paragraph{Open parity-alarm experiment.}
A separate 205-subtree, 5025-leaf source replaces emitted equality
$E(a,b)=4\ctP(a,b)+3$ by $\ctI(E(a,b),1)$, interprets raw
residue3 atoms as disequality, and tests cache parity with the four-leaf
scanner $(\mathrm{CD}_5,((\mathrm{CD}_5,\mathrm{CC}_2),\mathrm{CH}))$.
Constant assignments in an irreflexive model make both odd atom classes
false. However, the false witness $1=\operatorname{Rel}(0,0)$ adds
syntactic variable0 to the equality representation. Semantic validity
alone does not establish support-sensitive vacuity or the exact
formula/binder coupling in the gated bundle. A revised encoder and
reverse calibration remain to be proved; this experiment is not covered
by the theorem above.
\endgroup
